\documentclass[11pt,a4paper]{article}

% === 页面与字体 ===
\usepackage[margin=2.2cm]{geometry}
\usepackage{fontspec}
\setmainfont{DejaVu Serif}
\usepackage{xeCJK}
\setCJKmainfont{Droid Sans Fallback}
\setCJKfamilyfont{hei}{Droid Sans Fallback}
\newcommand{\hei}{\CJKfamily{hei}\bfseries}

% === 表格 ===
\usepackage{booktabs}
\usepackage{tabularx}
\usepackage{multirow}
\usepackage{array}
\usepackage{dcolumn}
\newcolumntype{d}[1]{D{.}{.}{#1}}
\usepackage{placeins}

% === caption 中文格式 ===
\usepackage{caption}
\renewcommand{\tablename}{表}
\captionsetup[table]{labelsep=quad, font={bf}, justification=centering}

% === 其他 ===
\usepackage{amsmath}
\usepackage{hyperref}
\hypersetup{colorlinks=false}
\usepackage{setspace}
\onehalfspacing

\usepackage{fancyhdr}
\pagestyle{fancy}
\fancyhf{}
\rhead{\small 数据要素利用与资本市场定价效率}
\lhead{\small 实证结果汇总}
\cfoot{\thepage}

\title{{\hei\Large 数据要素利用与资本市场定价效率} \\[8pt]
{\large 研究设计与实证结果}}
\author{}
\date{2026年3月}

\begin{document}
\maketitle
\thispagestyle{fancy}

% ============================================================
%                    研究设计
% ============================================================

\section*{三、研究设计}

\subsection*{（一）数据来源}

2011年国务院发布《关于深化科技体制改革加快国家创新体系建设的意见》标志着数据要素的战略价值开始进入政策视野，此后《促进大数据发展行动纲要》（2015年）、"数据二十条"（2022年）以及《企业数据资源相关会计处理暂行规定》（2024年）等一系列政策相继出台，推动了数据要素从概念萌芽到制度化确认的完整演进。本文即以2011年为样本起点，选取2011---2024年我国沪深A股上市公司为研究对象，在剔除金融及保险类（证监会J类）、ST以及*ST处理、上市不足一年以及变量缺失的数据后得到43\,847个年度-企业观测值，涵盖5\,178家独立企业。为了减轻极端值的影响，本文对所有连续变量进行上下1\%的缩尾处理。文中的数据要素利用水平指标采用文本分析方法从上市公司年度报告全文中提取得到，原始年报文件来源于上市公司向证监会提交的法定年报TXT格式文件；因变量股价延迟的构造基于股票日收益率数据，其他财务数据、治理数据和市场数据来自于CSMAR公司研究数据库。

\subsection*{（二）变量定义}

\noindent\textbf{1.\hspace{0.5em}被解释变量：资本市场定价效率（\textit{PriceDelay}）}

\vspace{2pt}

本文以股价延迟衡量资本市场定价效率。股价延迟反映企业特质信息从产生到融入股价之间的时滞，值越大表示信息反映越慢、定价效率越低。借鉴Hou和Moskowitz（2005）的做法，对企业$i$在年份$t$的日收益率序列（至少120个交易日），分别估计受限模型和无限制模型：

受限模型：
\vspace{-6pt}
\[ r_{i,d} = \alpha + \beta_0 \cdot r_{m,d} + \varepsilon_{i,d} \]
\vspace{-6pt}
无限制模型：
\vspace{-6pt}
\[ r_{i,d} = \alpha + \beta_0 \cdot r_{m,d} + \sum_{k=1}^{5}\beta_k \cdot r_{m,d-k} + \varepsilon_{i,d} \]
\vspace{-6pt}

其中$r_{i,d}$为企业$i$在第$d$日的个股收益率，$r_{m,d}$为沪深A股市值加权市场收益率。股价延迟定义为：
\vspace{-6pt}
\[ PriceDelay_{i,t} = 1 - \frac{R^2_{\text{restricted}}}{R^2_{\text{unrestricted}}} \]
\vspace{-6pt}

取值范围为$[0,\,1]$，值越大表示股价对市场信息的反应越滞后，定价效率越低。相比股价同步性（Morck等，2000）等常用指标，股价延迟更直接地捕捉信息向价格的转换速度而非信息结构的分解。在稳健性检验中，本文同时采用股价同步性（\textit{SYNCH}）作为替代度量。

\vspace{6pt}
\noindent\textbf{2.\hspace{0.5em}解释变量：数据要素利用水平（\textit{DU\_kw}）}

\vspace{2pt}

数据资产及其相关活动长期处于会计体系之外，难以通过传统财务指标直接观察，文本分析因此成为获取数据要素利用信息的重要途径（李世刚等，2025）。本文通过构建数据要素利用关键词库，沿数据价值链的不同环节设计了五维度关键词体系：数据存量（15个关键词，如大数据、数据库、数据中心等）、数据开发能力（19个关键词，如数据挖掘、机器学习、数字化转型等）、数据驱动应用（19个关键词，如精准营销、智能推荐、风控模型等）、数据价值变现（11个关键词，如数据资产、数据交易、数据入表等）和数据治理（9个关键词，如数据安全、数据隐私、数据合规等），共计73个关键词。这一分类方案借鉴了CNRDS数字化转型指标的理论框架，同时针对朱康和汤甬（2025）采用的两维度体系进行了细化，新增数据开发能力、数据价值变现和数据治理三个维度，更加全面地覆盖企业数据要素利用的全生命周期。然后采用文本分析方法在年报全文中进行长词优先匹配，统计各企业的数据要素利用词频。核心指标定义为：
\vspace{-6pt}
\[ DU_{kw} = \frac{\text{关键词总出现次数}}{\text{年报全文总字数}} \times 10000 \]
\vspace{-6pt}

即每万字中数据相关关键词的出现次数。为提高测度的稳健性，在主指标\textit{DU\_kw}的基础上，本文构造了两个备选指标：$DU\_kw\_ln = \ln(1 + DU_{kw})$，即对频率进行对数变换以处理右偏分布；$DU\_sub\_ln = \ln(1 + DU_{sub})$，其中$DU_{sub}$为实质利用次数，其构造基于对关键词出现上下文的判断------关键词周围100字窗口内同时包含具体数字、百分比表述或动作词与对象搭配时，判定为实质利用，否则视为概念性提及。完整的关键词清单列示于附录。

\vspace{6pt}
\noindent\textbf{3.\hspace{0.5em}控制变量}

\vspace{2pt}

参考已有关于股价延迟影响因素的相关研究（Hou和Moskowitz，2005；侯青川等，2016），本文从企业财务表现、治理结构、市场流动性和信息环境四个方面选取了控制变量：企业规模\textit{Size}=ln(总资产)、财务杠杆\textit{Lev}=总负债/总资产、总资产收益率\textit{ROA}=净利润/总资产、托宾Q值\textit{TobinQ}=(股权市值+负债账面值)/总资产、上市年限\textit{Age}=ln(当年-上市年份+1)、营收增长率\textit{Growth}、董事会规模\textit{BoardSize}=ln(董事会人数)、独立董事占比\textit{IndepRatio}=独立董事人数/董事会总人数、两职合一\textit{Dual}、第一大股东持股比例\textit{Top1Share}、产权性质\textit{SOE}、机构投资者持股比例\textit{InstHold}、Amihud非流动性指标\textit{Amihud}、分析师覆盖\textit{Analyst}=ln(1+跟踪分析师人数)以及审计意见\textit{AuditType}（标准无保留意见取1）。

\subsection*{（三）模型设计}

本文构建回归模型（1）识别数据要素利用对资本市场定价效率的影响效应：

\begin{equation}
PriceDelay_{i,t} = \alpha_0 + \alpha_1 DU_{i,t} + \alpha_2 Control_{i,t} + \mu_i + \lambda_t + \varepsilon_{i,t} \tag{1}
\end{equation}

其中，$i$代表企业，$t$代表年份，$PriceDelay$表示股价延迟，$DU$表示企业数据要素利用水平，$Control$表示控制变量，$\mu_i$和$\lambda_t$分别表示个体（Firm）及年份（Year）固定效应，$\varepsilon_{i,t}$表示随机误差项。为了避免行业内企业相关性和年份内系统性冲击对推断的影响，标准误在行业$\times$年份层面进行聚类调整。若假说H1a成立，即信息供给效应占主导，则$\alpha_1$应显著为负。

为检验数据要素利用影响定价效率的传导渠道，本文进一步构建如下路径模型：

\begin{equation}
Med_{i,t} = \beta_0 + \beta_1 DU_{i,t} + \beta_2 Control_{i,t} + \mu_i + \lambda_t + \varepsilon_{i,t} \tag{2}
\end{equation}

其中，$Med$分别为分析师覆盖（\textit{Analyst}）、金融资产占比（\textit{FinAsset}）和机构持股比例（\textit{InstHold}）。采用路径系数的直接呈现来识别传导渠道的存在性，而非传统的Sobel检验或Bootstrap分解（江艇，2022）。

\vspace{8pt}

% ============================================================
\FloatBarrier
% ============================================================
%                        表 1  描述性统计
% ============================================================

\begin{table}[htbp]
\centering
\caption{变量的描述性统计}
\label{tab:desc}
\small
\begin{tabular}{l*{7}{c}}
\toprule
变量 & 观测值 & 均值 & 标准差 & 最小值 & P25 & 中位数 & 最大值 \\
\midrule
\multicolumn{8}{l}{\textit{被解释变量}} \\
\textit{PriceDelay} & 43\,847 & 0.110 & 0.120 & 0.005 & 0.033 & 0.068 & 0.664 \\[4pt]
\multicolumn{8}{l}{\textit{核心解释变量}} \\
\textit{DU\_kw}      & 43\,847 & 1.070 & 1.695 & 0.000 & 0.193 & 0.478 & 10.060 \\
\textit{DU\_kw\_ln}  & 43\,847 & 2.415 & 1.250 & 0.000 & 1.609 & 2.398 & 5.481 \\
\textit{DU\_sub}     & 43\,847 & 12.279 & 24.132 & 0.000 & 1.000 & 4.000 & 152.000 \\
\textit{DU\_ratio}   & 43\,847 & 0.621 & 0.379 & 0.000 & 0.333 & 0.750 & 1.000 \\[4pt]
\multicolumn{8}{l}{\textit{控制变量}} \\
\textit{Size}        & 43\,847 & 22.643 & 0.972 & 20.949 & 21.934 & 22.481 & 25.667 \\
\textit{Lev}         & 43\,847 & 0.434 & 0.205 & 0.055 & 0.271 & 0.427 & 0.923 \\
\textit{ROA}         & 43\,847 & 0.027 & 0.068 & $-$0.279 & 0.009 & 0.031 & 0.191 \\
\textit{TobinQ}      & 43\,847 & 2.027 & 1.314 & 0.830 & 1.219 & 1.604 & 8.443 \\
\textit{Age}         & 43\,847 & 2.282 & 0.635 & 1.099 & 1.792 & 2.398 & 3.219 \\
\textit{Growth}      & 43\,847 & 0.127 & 0.371 & $-$0.591 & $-$0.056 & 0.075 & 2.157 \\
\textit{BoardSize}   & 43\,847 & 2.109 & 0.198 & 1.609 & 1.946 & 2.197 & 2.639 \\
\textit{IndepRatio}  & 43\,847 & 0.379 & 0.054 & 0.333 & 0.333 & 0.364 & 0.571 \\
\textit{Dual}        & 43\,847 & 0.287 & 0.453 & 0.000 & 0.000 & 0.000 & 1.000 \\
\textit{Top1Share}   & 43\,847 & 33.034 & 14.756 & 8.110 & 21.600 & 30.550 & 74.020 \\
\textit{SOE}         & 43\,847 & 0.340 & 0.474 & 0.000 & 0.000 & 0.000 & 1.000 \\
\textit{InstHold}    & 43\,847 & 42.081 & 24.331 & 0.350 & 21.540 & 43.372 & 90.708 \\
\textit{Amihud}      & 43\,847 & 0.042 & 0.043 & 0.002 & 0.015 & 0.030 & 0.302 \\
\textit{Analyst}     & 43\,847 & 2.489 & 1.739 & 0.000 & 0.693 & 2.639 & 5.656 \\
\textit{AuditType}   & 43\,847 & 0.963 & 0.188 & 0.000 & 1.000 & 1.000 & 1.000 \\
\bottomrule
\end{tabular}

\vspace{4pt}
{\footnotesize 注：连续变量均在1\%和99\%分位数处缩尾处理。\textit{Size}=ln(总资产)，\textit{Age}=ln(上市年限)，\textit{Analyst}=ln(1+分析师人数)，\textit{BoardSize}=ln(董事会人数)。样本为2011--2024年沪深A股非金融非ST上市公司。}
\end{table}


% ============================================================
\FloatBarrier
% ============================================================
%                        表 2  基准回归
% ============================================================

\begin{table}[htbp]
\centering
\caption{基准回归检验}
\label{tab:baseline}
\small
\setlength{\tabcolsep}{5pt}
\begin{tabular}{l*{6}{c}}
\toprule
 & (1) & (2) & (3) & (4) & (5) & (6) \\
 & \textit{PriceDelay} & \textit{PriceDelay} & \textit{PriceDelay} & \textit{PriceDelay} & \textit{PriceDelay} & \textit{PriceDelay} \\
\midrule
\textit{DU\_kw}
 & $-$0.0053$^{***}$ & $-$0.0042$^{***}$ & & & $-$0.0018$^{***}$ & \\
 & (0.0010) & (0.0010) & & & (0.0006) & \\[4pt]
\textit{DU\_kw\_ln}
 & & & $-$0.0041$^{***}$ & & & $-$0.0018$^{***}$ \\
 & & & (0.0012) & & & (0.0007) \\[4pt]
\textit{DU\_sub\_ln}
 & & & & $-$0.0035$^{***}$ & & \\
 & & & & (0.0007) & & \\[4pt]
\textit{Size}
 & & 0.020$^{***}$ & & & & \\
 & & (0.005) & & & & \\[2pt]
\textit{Lev}
 & & 0.023$^{***}$ & & & & \\
 & & (0.007) & & & & \\[2pt]
\textit{ROA}
 & & $-$0.037$^{**}$ & & & & \\
 & & (0.015) & & & & \\[2pt]
\textit{TobinQ}
 & & 0.010$^{***}$ & & & & \\
 & & (0.001) & & & & \\[2pt]
\textit{Age}
 & & $-$0.026$^{***}$ & & & & \\
 & & (0.007) & & & & \\[2pt]
\textit{Growth}
 & & 0.010$^{***}$ & & & & \\
 & & (0.002) & & & & \\[2pt]
\textit{BoardSize}
 & & $-$0.014$^{**}$ & & & & \\
 & & (0.005) & & & & \\[2pt]
\textit{IndepRatio}
 & & $-$0.036$^{**}$ & & & & \\
 & & (0.018) & & & & \\[2pt]
\textit{Dual}
 & & $-$0.002 & & & & \\
 & & (0.002) & & & & \\[2pt]
\textit{Top1Share}
 & & $-$0.000 & & & & \\
 & & (0.000) & & & & \\[2pt]
\textit{SOE}
 & & 0.000 & & & & \\
 & & (0.004) & & & & \\[2pt]
\textit{InstHold}
 & & 0.000$^{***}$ & & & & \\
 & & (0.000) & & & & \\[2pt]
\textit{Amihud}
 & & 0.146$^{***}$ & & & & \\
 & & (0.023) & & & & \\[2pt]
\textit{Analyst}
 & & $-$0.010$^{***}$ & & & & \\
 & & (0.001) & & & & \\[2pt]
\textit{AuditType}
 & & $-$0.014$^{***}$ & & & & \\
 & & (0.004) & & & & \\[4pt]
\midrule
控制变量 & 否 & 是 & 是 & 是 & 是 & 是 \\
固定效应 & 公司+年份 & 公司+年份 & 公司+年份 & 公司+年份 & 行业+年份 & 行业+年份 \\
$N$ & 43\,847 & 43\,847 & 43\,847 & 43\,847 & 34\,558 & 34\,558 \\
$R^2$ & & 0.425 & & & & \\
$R^2_{\text{within}}$ & & 0.029 & & & & \\
\bottomrule
\end{tabular}

\vspace{4pt}
{\footnotesize 注：括号内为标准误，聚类于行业$\times$年份层面。$^{***}$、$^{**}$、$^{*}$分别表示在1\%、5\%、10\%的水平上显著，下同。模型(2)为主报告模型，控制变量同时适用于模型(3)--(6)但限于篇幅未报告。经济意义：\textit{DU\_kw}增加一个标准差(1.70)$\to$ \textit{PriceDelay}降低约6.4\%。}
\end{table}


% ============================================================
\FloatBarrier
% ============================================================
%                    表 3  工具变量估计
% ============================================================

\begin{table}[htbp]
\centering
\caption{工具变量估计}
\label{tab:iv}
\small

\vspace{4pt}
{\hei Panel A：主结果（$N=38\,598$）}
\vspace{4pt}

\begin{tabular}{l*{3}{c}}
\toprule
 & (1) OLS & (2) 一阶段 & (3) 2SLS \\
 & \textit{PriceDelay} & \textit{DU\_kw} & \textit{PriceDelay} \\
\midrule
\textit{DU\_kw}
 & $-$0.0036$^{**}$ & & $-$0.0446$^{***}$ \\
 & (0.0014) & & (0.0161) \\[4pt]
\textit{IV1} (Toi\_s$\times$year\_c)
 & & 0.0008$^{***}$ & \\
 & & (0.0001) & \\[4pt]
\midrule
控制变量 & 是 & 是 & 是 \\
固定效应 & 公司+年份 & 公司+年份 & 公司+年份 \\
一阶段$F$ & & 92.2 & \\
\bottomrule
\end{tabular}

\vspace{14pt}
{\hei Panel B：替代口径}
\vspace{4pt}

\begin{tabular}{lccc}
\toprule
IV口径 & 一阶段$F$ & 2SLS系数 & 标准误 \\
\midrule
Toi\_s（综合指数，主） & 92.2 & $-$0.045$^{***}$ & (0.016) \\
Coi\_s（商用指数） & 59.8 & $-$0.066$^{***}$ & (0.019) \\
Goi\_s（政用指数） & 97.6 & $-$0.034$^{**}$ & (0.016) \\
\bottomrule
\end{tabular}

\vspace{14pt}
{\hei Panel C：Oster (2019) $\delta^*$界}
\vspace{4pt}

\begin{tabular}{lcc}
\toprule
固定效应规格 & $\delta^*$ ($R_{\max}=1.3\tilde{R}^2$) & 结论 \\
\midrule
公司+年份 & 9.03 & 稳健 \\
公司+行业$\times$年份 & 3.65 & 稳健 \\
\bottomrule
\end{tabular}

\vspace{6pt}
{\footnotesize 注：Panel A中IV1为2016年省级大数据发展指数(Toi\_s)乘以时间趋势(year$-$2016)。一阶段$F$远超Stock--Yogo临界值10。2SLS系数为OLS的12.2倍，可归因于\textit{DU\_kw}作为关键词计数的测量误差导致OLS存在衰减偏误，IV修正后效应更大。Panel B三个替代口径方向完全一致、全部显著。Panel C两个规格$\delta^*>1$，遗漏变量需达到可观测控制变量的3.7--9.0倍才能推翻结果。过度识别检验（IV1+IV2, Hansen $J=21.44$, $p<0.001$）拒绝联合外生性，故仅以IV1作为主报告。}
\end{table}


% ============================================================
\FloatBarrier
% ============================================================
%                    表 4  机制检验
% ============================================================

\begin{table}[htbp]
\centering
\caption{作用机制检验}
\label{tab:mechanism}
\small
\begin{tabular}{l*{3}{c}}
\toprule
 & (1) & (2) & (3) \\
 & \textit{Analyst} & \textit{FinAsset} & \textit{InstHold} \\
 & 分析师跟踪 & 金融资产占比 & 机构持股比例 \\
\midrule
\textit{DU\_kw}
 & 0.0255$^{***}$ & 0.0008$^{***}$ & $-$0.2518$^{***}$ \\
 & (0.0097) & (0.0003) & (0.0789) \\[6pt]
\midrule
控制变量 & 是 & 是 & 是 \\
固定效应 & 公司+年份 & 公司+年份 & 公司+年份 \\
$N$ & 43\,847 & 43\,847 & 43\,847 \\
$R^2$ & 0.794 & 0.671 & 0.916 \\
\bottomrule
\end{tabular}

\vspace{4pt}
{\footnotesize 注：分析师渠道是核心传导路径：数据利用高的企业吸引更多分析师跟踪，加速信息向股价传导。\textit{InstHold}方向为负，可解释为数据驱动型企业信息复杂度高，部分机构投资者选择回避，但分析师渠道的补偿效应更强。}
\end{table}


% ============================================================
\FloatBarrier
% ============================================================
%                    表 5  异质性分析
% ============================================================

\begin{table}[htbp]
\centering
\caption{异质性检验}
\label{tab:hetero}
\footnotesize
\setlength{\tabcolsep}{3.5pt}
\begin{tabular}{l*{8}{c}}
\toprule
 & \multicolumn{2}{c}{产权性质} & \multicolumn{2}{c}{企业规模} & \multicolumn{2}{c}{分析师覆盖} & \multicolumn{2}{c}{行业属性} \\
\cmidrule(lr){2-3}\cmidrule(lr){4-5}\cmidrule(lr){6-7}\cmidrule(lr){8-9}
 & (1) & (2) & (3) & (4) & (5) & (6) & (7) & (8) \\
 & 国企 & 民企 & 大企业 & 小企业 & 高覆盖 & 低覆盖 & 高科技 & 传统 \\
\midrule
\textit{DU\_kw}
 & $-$0.0060$^{***}$ & $-$0.0034$^{***}$ & $-$0.0033 & $-$0.0048$^{***}$ & $-$0.0055$^{***}$ & $-$0.0033$^{**}$ & $-$0.0044$^{***}$ & $-$0.0045$^{***}$ \\
 & (0.0015) & (0.0009) & (0.0022) & (0.0010) & (0.0010) & (0.0013) & (0.0014) & (0.0015) \\[6pt]
\midrule
控制变量 & 是 & 是 & 是 & 是 & 是 & 是 & 是 & 是 \\
固定效应 & \multicolumn{8}{c}{公司+年份} \\
$N$ & 14\,414 & 29\,371 & 18\,885 & 18\,895 & 21\,838 & 21\,459 & 9\,525 & 34\,050 \\
\bottomrule
\end{tabular}

\vspace{4pt}
{\footnotesize 注：大/小企业以\textit{Size}年度中位数划分；高/低分析师覆盖以\textit{Analyst}年度中位数划分；高科技行业包括计算机(C39)、软件(I65)、专用设备(C35)等。小企业效应强于大企业（信息不对称更严重），高分析师覆盖组效应强度约为低覆盖组的1.7倍（与分析师中介机制一致）。所有子样本DU\_kw系数均显著为负。}
\end{table}


% ============================================================
\FloatBarrier
% ============================================================
%                    表 6  稳健性检验
% ============================================================

\begin{table}[htbp]
\centering
\caption{稳健性检验}
\label{tab:robust}
\footnotesize
\setlength{\tabcolsep}{4pt}
\begin{tabular}{l*{6}{c}}
\toprule
 & (1) & (2) & (3) & (4) & (5) & (6) \\
 & \multicolumn{3}{c}{替换被解释变量} & PSM & 剔除 & 仅主板 \\
\cmidrule(lr){2-4}
 & \textit{SYNCH} & \textit{SYNCH\_ln} & \textit{SYNCH} & \textit{PriceDelay} & \textit{PriceDelay} & \textit{PriceDelay} \\
 & Firm+Year & Firm+Year & Ind+Year & Firm+Year & Firm+Year & Firm+Year \\
\midrule
\textit{DU\_kw}
 & 0.0110$^{*}$ & & 0.0111$^{**}$ & $-$0.0048$^{***}$ & $-$0.0050$^{***}$ & $-$0.0039$^{**}$ \\
 & (0.0066) & & (0.0048) & (0.0013) & (0.0010) & (0.0016) \\[4pt]
\textit{DU\_kw\_ln}
 & & 0.0152$^{**}$ & & & & \\
 & & (0.0062) & & & & \\[4pt]
\midrule
控制变量 & 是 & 是 & 是 & 是 & 是 & 是 \\
固定效应 & 公司+年份 & 公司+年份 & 行业+年份 & 公司+年份 & 公司+年份 & 公司+年份 \\
$N$ & 29\,444 & 29\,444 & 29\,924 & 31\,381 & 33\,878 & 29\,666 \\
\bottomrule
\end{tabular}

\vspace{4pt}
{\footnotesize 注：模型(1)--(3)替换被解释变量为股价同步性（\textit{SYNCH}），系数方向为正（数据利用提高同步性），与\textit{PriceDelay}方向相反，符合预期。模型(4)为PSM 1:1最近邻匹配（caliper=0.05），平衡性良好（14/14变量SMD$<$12\%）。模型(5)剔除2024年样本。模型(6)仅保留主板样本。}
\end{table}


% ============================================================
\FloatBarrier
% ============================================================
%                    表 7  投资组合检验
% ============================================================

\begin{table}[htbp]
\centering
\caption{基于数据利用分组的投资组合$\alpha$检验}
\label{tab:portfolio}
\small

\vspace{4pt}
{\hei Panel A：各分位组月度$\alpha$（$N=156$个月）}
\vspace{4pt}

\begin{tabular}{lcccc}
\toprule
组合 & CAPM & FF3 & FF5 & FF5+MOM \\
\midrule
Q1（最低利用）
 & 0.0025 & $-$0.0029$^{**}$ & $-$0.0053$^{***}$ & $-$0.0064$^{***}$ \\
 & (0.81) & ($-$2.08) & ($-$3.15) & ($-$3.81) \\[3pt]
Q2
 & 0.0027 & $-$0.0017 & $-$0.0047$^{***}$ & $-$0.0055$^{***}$ \\
 & (0.84) & ($-$1.58) & ($-$3.11) & ($-$3.64) \\[3pt]
Q3
 & 0.0037 & $-$0.0004 & $-$0.0039$^{***}$ & $-$0.0044$^{***}$ \\
 & (1.16) & ($-$0.43) & ($-$3.06) & ($-$3.40) \\[3pt]
Q4
 & 0.0045 & 0.0008 & $-$0.0031$^{**}$ & $-$0.0033$^{**}$ \\
 & (1.23) & (0.81) & ($-$2.20) & ($-$2.31) \\[3pt]
Q5（最高利用）
 & 0.0054 & 0.0009 & $-$0.0039 & $-$0.0041 \\
 & (1.04) & (0.38) & ($-$1.58) & ($-$1.60) \\[3pt]
\midrule
DAT (Q5$-$Q1)
 & 0.0029 & 0.0037 & 0.0014 & 0.0023 \\
 & (0.77) & (1.17) & (0.44) & (0.71) \\
\bottomrule
\end{tabular}

\vspace{14pt}
{\hei Panel B：入表前后对比}
\vspace{4pt}

\begin{tabular}{lcc}
\toprule
期间 & DAT FF3 $\alpha$ & $t$值 \\
\midrule
入表前（2012--2023） & $+$0.21\% & 0.63 \\
入表后（2024--） & $+$2.08\% & 2.24$^{**}$ \\
\bottomrule
\end{tabular}

\vspace{6pt}
{\footnotesize 注：每年6月按\textit{DU\_kw}五分组，7月至次年6月持有。Panel A括号内为$t$统计量。Q1（低数据利用）在FF3模型下$\alpha$显著为负（月均$-$0.29\%，年化约$-$3.4\%），说明市场对数据利用不足的企业系统性高估。Panel B中，2024年入表新规后高低组收益差异从不显著跳升至$+$2.08\%/月（$t=2.24$），说明政策使数据利用差异更为可见。GRS检验$F=1.44$（$p=0.213$），不拒绝所有$\alpha$联合为零。}
\end{table}


% ============================================================
\FloatBarrier
% ============================================================
%                    附表 A1  DID 检验
% ============================================================

\begin{table}[htbp]
\centering
\caption*{{\bf 附表A1\quad DID检验（2024年数据资产入表新规）}}
\label{tab:did}
\small

\vspace{4pt}
{\hei Panel A：主结果}
\vspace{4pt}

\begin{tabular}{l*{5}{c}}
\toprule
 & (1) & (2) & (3) & (4) & (5) \\
 & \multicolumn{3}{c}{二值DID} & \multicolumn{2}{c}{连续DID} \\
\cmidrule(lr){2-4}\cmidrule(lr){5-6}
处理组定义 & 2017--21 & 2011--21 & 旧分组 & 2019--24 & 2021--24 \\
 & 均值 & 均值 & 2019--23 & 强度 & 强度 \\
\midrule
\textit{Treat}$\times$\textit{Post}
 & 0.0196$^{***}$ & 0.0201$^{***}$ & 0.0075 & 0.0097$^{***}$ & 0.0087$^{***}$ \\
 & (0.0042) & (0.0035) & (0.0054) & (0.0018) & (0.0021) \\[4pt]
\midrule
$N$ & 16\,005 & 16\,005 & 16\,005 & 22\,471 & 16\,005 \\
\bottomrule
\end{tabular}

\vspace{14pt}
{\hei Panel B：安慰剂检验（虚假政策年份）}
\vspace{4pt}

\begin{tabular}{lcccc}
\toprule
虚假年份 & 系数 & 标准误 & $t$值 & $p$值 \\
\midrule
2019 & $-$0.0114 & (0.0048) & $-$2.39 & 0.017$^{**}$ \\
2020 & $-$0.0127 & (0.0058) & $-$2.21 & 0.028$^{**}$ \\
2021 & $-$0.0062 & (0.0051) & $-$1.20 & 0.233 \\
2022 & 0.0078 & (0.0044) & 1.77 & 0.078$^{*}$ \\
\textbf{2024（真实）} & \textbf{0.0196} & \textbf{(0.0042)} & \textbf{4.68} & $<$\textbf{0.001}$^{***}$ \\
\bottomrule
\end{tabular}

\vspace{14pt}
{\hei Panel C：平行趋势检验（基期=2023）}
\vspace{4pt}

\begin{tabular}{lcc}
\toprule
年份交互 & 系数 & $t$值 \\
\midrule
\textit{Treat}$\times$2019 & $-$0.0056 & $-$0.92 \\
\textit{Treat}$\times$2020 & $-$0.0101 & $-$1.23 \\
\textit{Treat}$\times$2021 & $-$0.0159$^{**}$ & $-$2.41 \\
\textit{Treat}$\times$2022 & 0.0008 & 0.19 \\
\textit{Treat}$\times$2024 & 0.0048$^{***}$ & 4.37 \\
\bottomrule
\end{tabular}

\vspace{6pt}
{\footnotesize 注：Panel A窗口为2021--2024，公司+年份固定效应，聚类于行业$\times$年份。Panel B中2019/2020虚假政策年安慰剂检验显著，表明处理组在政策实施前已系统性异于控制组，平行趋势假设不满足。入表新规为普遍性政策，缺乏真正的对照组；处理组定义（\textit{DU\_kw}中位数分组）具有内生性。基于上述原因，DID设计存在内部效度问题，仅作为附录参考。}
\end{table}


\end{document}
